\begin{table}[htbp]
\centering
\caption{主要模型迴歸結果}
\label{tab:reg}
\footnotesize
\setlength{\tabcolsep}{6pt}
\renewcommand{\arraystretch}{1.12}
\begin{tabular}{l rrrrrr}
\toprule
& 幸福感 & 自尊 & 生活滿意度 & 好友支持 & 孤獨感$^{\dagger}$ & 心理困擾$^{\dagger}$ \\
\midrule
\multicolumn{7}{l}{\textbf{A. 四象限分類模型}（對照組：雙低型）} \\
\quad 線上優勢型 & $-0.172^{***}$ & $-0.137^{***}$ & $-0.168^{***}$ & $0.205^{***}$ & $0.326^{***}$ & $0.287^{***}$ \\
 & (0.040) & (0.027) & (0.027) & (0.026) & (0.040) & (0.030) \\[1pt]
\quad 線下優勢型 & $0.720^{***}$ & $0.612^{***}$ & $0.452^{***}$ & $0.413^{***}$ & $-0.268^{***}$ & $-0.264^{***}$ \\
 & (0.041) & (0.026) & (0.025) & (0.026) & (0.031) & (0.024) \\[1pt]
\quad 雙高型 & $0.604^{***}$ & $0.547^{***}$ & $0.340^{***}$ & $0.474^{***}$ & $-0.109^{***}$ & $-0.128^{***}$ \\
 & (0.036) & (0.023) & (0.023) & (0.022) & (0.030) & (0.024) \\[1pt]
\quad $R^2$ & .113 & .193 & .116 & .118 & .054 & .083 \\
\midrule
\multicolumn{7}{l}{\textbf{B. 連續分數交互作用模型}} \\
\quad 線上自我認同 & $-0.156^{***}$ & $-0.132^{***}$ & $-0.152^{***}$ & $0.094^{***}$ & $0.236^{***}$ & $0.211^{***}$ \\
 & (0.017) & (0.011) & (0.011) & (0.012) & (0.016) & (0.012) \\[1pt]
\quad 線下自我認同 & $0.617^{***}$ & $0.520^{***}$ & $0.391^{***}$ & $0.266^{***}$ & $-0.309^{***}$ & $-0.280^{***}$ \\
 & (0.016) & (0.011) & (0.012) & (0.012) & (0.019) & (0.014) \\[1pt]
\quad 線上 $\times$ 線下 & $0.098^{***}$ & $0.085^{***}$ & $0.065^{***}$ & $-0.020$ & $-0.056^{***}$ & $-0.048^{***}$ \\
 & (0.016) & (0.012) & (0.012) & (0.012) & (0.017) & (0.013) \\[1pt]
\quad $R^2$ & .161 & .262 & .161 & .169 & .078 & .109 \\
\midrule
背景控制變數 & 是 & 是 & 是 & 是 & 是 & 是 \\
$N$ & 7,949 & 7,963 & 7,995 & 8,010 & 7,958 & 7,958 \\
\bottomrule
\end{tabular}

\vspace{0.4em}
\parbox{\linewidth}{\scriptsize 註：括號內為 HC3 穩健標準誤，$p$ 值依六項指標進行 FDR 校正；背景控制變數之係數未列出。模型 B 的線上與線下自我認同分數經平均數中心化。$^{*}p<.05$，$^{**}p<.01$，$^{***}p<.001$。$^{\dagger}$分數越高代表狀態越差。}
\end{table}
